% BHU PYQ -- Transcribed & Corrected
% Paper: ESMJ43 / ESMN41: Environmental Geology and Hydrogeology
% Exam: B.Sc. (Hons.) Semester IV Examination, 2025-26 (NEP)
% Department: Department of Geology

\documentclass[11pt,a4paper]{article}
\usepackage[a4paper,top=2.5cm,bottom=2.5cm,left=2.5cm,right=2.5cm,headheight=14pt]{geometry}
\usepackage{amsmath,amssymb}
\usepackage[T1]{fontenc}
\usepackage[utf8]{inputenc}
\usepackage{lmodern,microtype,fancyhdr,enumitem,hyperref,lastpage,parskip}

\hypersetup{colorlinks=true,urlcolor=black,linkcolor=black,pdftitle={ESMJ43: Environmental Geology & Hydrogeology -- BHU Sem IV 2025-26 (NEP)}}
\pagestyle{fancy}\fancyhf{}
\fancyhead[L]{\small   \textit{Banaras Hindu University}}
\fancyhead[R]{\small   \textit{ESMJ43 / ESMN41: Environmental Geology \& Hydrogeology}}
\fancyfoot[C]{\small Page  \thepage\ of \pageref{LastPage}}


\newlist{parts}{enumerate}{2}
\setlist[parts,1]{label=(\roman*),leftmargin=*,itemsep=4pt,topsep=2pt}


\begin{document}


\begin{center}
  {\large\bfseries BANARAS HINDU UNIVERSITY}\\[2pt]
  {
\normalsize B.Sc. (Hons.) Semester IV Examination 2025-26 (NEP)}\\[6pt]
  \rule{0.8\linewidth}{0.5pt}\\[6pt]
  {\large\bfseries DEPARTMENT OF GEOLOGY}\\[4pt]
  {\large\bfseries Paper Code: ESMJ43 / ESMN41 --- Environmental Geology and Hydrogeology}\\[4pt]
  {
\normalsize Time: 2:30 Hours\hfill Maximum Marks: 70}
\end{center}

\rule{\linewidth}{0.4pt}\smallskip



\noindent   \textit{Instruction: Write your Roll No. at the top immediately on the receipt of this question paper.}\\[2pt]



\noindent   \textit{Note: Answer five questions in all, selecting at least two questions each from Section--A and Section--B. Section--C is compulsory. All questions carry equal marks (14 marks each).}
\bigskip

\section*{SECTION -- A}




\noindent   \textbf{Question 1} \hfill [14 Marks]\\
Explain how global climate drivers, orbital parameters, and surface albedo feedback loops modulate global temperatures.

\bigskip




\noindent   \textbf{Question 2} \hfill [14 Marks]\\
What is Darcy's Law? Derive its equation and comment on the validity of Darcy's Law with respect to Reynolds number.

\bigskip




\noindent   \textbf{Question 3} \hfill [14 Marks]\\
What is the hydrological cycle? Explain briefly the various components of the Hydrological Cycle with a neat sketch.

\bigskip
\section*{SECTION -- B}




\noindent   \textbf{Question 4} \hfill [14 Marks]\\
Briefly discuss the subsurface methods of groundwater exploration.

\bigskip




\noindent   \textbf{Question 5} \hfill [14 Marks]\\
Describe the basic terms related to aquifer properties. Calculate the groundwater storage for an aquifer having an areal extent of $10 \text{ km}^2$ and a saturated thickness of $15 \text{ m}$. Consider the specific yield of the aquifer to be $20\%$.

\bigskip




\noindent   \textbf{Question 6} \hfill [14 Marks]\\
Attempt any   \textbf{two} of the following ($2  \times 7 = 14$ Marks):

\begin{parts}
  \item Hydro-chemical properties used to characterize groundwater quality
  \item Causes of groundwater depletion
  \item Geological mechanism of earthquake generation and seismic hazards
  \item Landslides and their causes
\end{parts}

\bigskip
\section*{SECTION -- C (Compulsory)}




\noindent   \textbf{Question 7} \hfill [14 Marks]\\
Answer the following in 2 to 3 sentences each ($7  \times 2 = 14$ Marks):

\begin{parts}
  \item Environmental legislation
  \item Role of Stratosphere
  \item Population growth
  \item Aquifer
  \item Water table and Water level
  \item Intrinsic Permeability
  \item Effluent and Influent conditions
\end{parts}

\end{document}
